\begin{figure}[t]
\centering
\begin{tikzpicture}[scale=.88,>=Latex,font=\small]
  \draw[thick] (0,.65) -- (5,.65);
  \draw[thick] (0,1.85) -- (5,1.85);
  \draw[->] (4.65,.65) -- (4.25,.65);
  \draw[->] (4.65,1.85) -- (4.25,1.85);
  \foreach \x/\lab in {1/1,2.15/2,4/N}{
    \draw[thick,->] (\x,-.05) -- (\x,2.55);
    \draw[fill=white] (\x-.12,.53) rectangle (\x+.12,.77);
    \draw[fill=white] (\x-.12,1.73) rectangle (\x+.12,1.97);
    \node[below] at (\x,-.05) {$\lab$};
  }
  \node at (3.05,1.25) {$\cdots$};
  \draw[thick] (0,.65) .. controls (-.55,.65) and (-.55,1.85) .. (0,1.85);
  \draw[thick] (5,.65) .. controls (5.55,.65) and (5.55,1.85) .. (5,1.85);
  \node[anchor=east,align=right] at (-.5,1.25) {$K_{e_\ell}(r)^T$};
  \node[anchor=west] at (5.5,1.25) {$K_{e_r}(t)$};
  \node[above] at (3.25,1.88) {$A:\ t$};
  \node[below] at (2.5,-.85) {$B:\ r=b-t$};
  \node[above] at (2.5,2.7) {Double-row contraction};
  \begin{scope}[shift={(8.3,.35)},scale=1.15]
    \coordinate (O) at (0,0);
    \coordinate (V) at (1.5,0);
    \coordinate (U) at (.75,1.299);
    \coordinate (UV) at (2.25,1.299);
    \coordinate (C1) at (.75,.433);
    \coordinate (C2) at (1.5,.866);
    \draw[gray] (O)--(V)--(UV)--(U)--cycle;
    \draw[gray] (U)--(V);
    \draw[thick] (.375,.6495)--(C1)--(.75,0);
    \draw[thick] (C1)--(C2);
    \draw[thick] (1.5,1.299)--(C2)--(1.875,.6495);
    \fill (C1) circle (1.6pt);
    \fill (C2) circle (1.6pt);
    \node[left] at (.375,.6495) {$w$};
    \node[below] at (.75,0) {$s$};
    \node[above] at (1.5,1.299) {$n$};
    \node[right] at (1.875,.6495) {$e$};
    \node[above] at (1.125,1.6) {Physical lozenge};
  \end{scope}
\end{tikzpicture}
\caption{The horizontal and vertical lines denote spin spaces, and the boundary curves denote the displayed tensors. The auxiliary row $B$ acts first on the site spins, then row $A$; products act from right to left. The two auxiliary ends are contracted with the displayed boundary tensors. This is a tensor schematic, whereas the inset shows the actual pair of equilateral triangles at the physical specialization. The black segments in the inset are the available honeycomb links, not one path configuration. The two dots are the centers counted by the path length: an $s$--$w$ or $e$--$n$ arc visits one center, and a single arc between the two pairs visits both.}
\label{fig:strip-double-row}
\end{figure}
